\documentclass[11pt]{article}
\usepackage[margin=1in]{geometry}
\usepackage{amsmath,amssymb,amsthm,booktabs,hyperref}
\newtheorem{theorem}{Theorem}
\newtheorem{proposition}[theorem]{Proposition}
\newtheorem{lemma}[theorem]{Lemma}
\theoremstyle{remark}\newtheorem{remark}[theorem]{Remark}
\newcommand{\dx}{\,dx}
\newcommand{\nrm}[1]{\lVert #1\rVert}
\newcommand{\Amin}{a_{\min}}
\title{Lemma O: explicit short-time stability for the graph form of curve diffusion flow}
\author{}\date{}
\begin{document}\maketitle

\section{Statement}
Let $h(x,t)$ be smooth on $[-2,2]\times[0,t_1]$ and solve
\begin{equation}\label{eq:graph}
h_t=-\partial_x G(h),\qquad G(h)=a(h_x)h_{xxx}+b(h_x)h_{xx}^2,\quad a(p)=(1+p^2)^{-2},\ b(p)=-3p(1+p^2)^{-3}.
\end{equation}
(\emph{Derivation, checked symbolically in \texttt{sanity.py}:} for a graph, $\kappa=h_{xx}(1+h_x^2)^{-3/2}$, $\partial_s=(1+h_x^2)^{-1/2}\partial_x$, and the normal velocity $V=-\kappa_{ss}$ gives $h_t=\sqrt{1+h_x^2}\,V=-\partial_x\big[(1+h_x^2)^{-1/2}\kappa_x\big]$; expanding $\kappa_x=h_{xxx}(1+h_x^2)^{-3/2}-3h_xh_{xx}^2(1+h_x^2)^{-5/2}$ gives \eqref{eq:graph}.)
Write $h_0=h(\cdot,0)$, $g=h-h_0$, $g_j=\partial_x^jg$, $I_0=(-2,2)$, $\nrm{\cdot}=L^2(I_0)$.

\paragraph{Hypotheses.} For all $t\in[0,t_1]$:
\begin{itemize}
\item[(A)] $|h_x|\le\frac14$, $|h_{xx}|\le K_2$ on $[-2,2]$, and $\int_{-2}^2h_{xxx}^2\le K_3^2$.
\item[(I)] $|h_0^{(j)}|\le M$ on $[-2,2]$ for $1\le j\le4$, and $\int_{-2}^2(h_0^{(5)})^2\le M^2$.
\end{itemize}

\begin{theorem}[Lemma O]\label{thm:O}
There are explicit $C_0,\dots,C_3$ (formulas \eqref{eq:consts} below, evaluated by \texttt{lemmaO\_consts.py}), depending only on $K_2,K_3,M$ (and, for $C_0$ only, additively on $\sqrt{t_1}$), such that for every $t\in[0,t_1]$ and $x\in[-1,1]$
\[
|h-h_0|\le C_0t^{1/2},\quad |h_x-h_{0,x}|\le C_1t^{1/2},\quad |h_{xx}-h_{0,xx}|\le C_2t^{1/2},\quad \int_{-1}^1(h_{xxx}-h_{0,xxx})^2\le C_3\,t .
\]
The proof needs \emph{no} smallness of $t_1$: the a priori bounds (A) are used as hypotheses on the whole interval, so $\delta_0(K_2,K_3,M)=+\infty$ formally. The statement is quantitatively informative only when $C_j\sqrt t$ is small compared with the trivial sizes given by (A),(I); see \S\ref{sec:use}, where the answer for the two use-cases is that the constants are \emph{too large} to be useful at the requested times.
\end{theorem}

\section{Preliminaries}
\paragraph{Consequences of (A),(I).} On $[-2,2]$, for $t\le t_1$:
\begin{equation}\label{eq:P}
|g_1|\le p_1:=\tfrac14+M,\quad |g_2|\le k_2:=K_2+M,\quad \nrm{g_3}\le\kappa:=K_3+2M,\quad \nrm{G(h)}\le \Gamma:=K_3+2B_0K_2^2 .
\end{equation}
($\nrm{h_0'''}\le 2M$ since $|I_0|^{1/2}=2$; $\nrm{b\,h_{xx}^2}\le B_0K_2^2\cdot2$, $|a|\le1$.)

\begin{lemma}[coefficients]\label{lem:coef}
Let $r=(16/17)^3$, $s=1+p^2$, $|p|\le\frac14$. Then $a_{\min}:=\min a=(17/16)^{-2}=0.8858\le a\le1$ and
\[
a'=-4ps^{-3},\ b=-3ps^{-3},\ b'=-3(1-5p^2)s^{-4},\ c_1:=2a'+2b=-14ps^{-3},\ c_2:=a'+2b=-10ps^{-3},
\]
\[ c_3:=a''+5b'=19(5p^2-1)s^{-4},\qquad c_4:=b''=-18p(5p^2-3)s^{-5},
\]
with $A_1:=\sup|a'|=r$, $B_0:=\sup|b|=\frac34r$, $B_1:=\sup|b'|=3$, $C_1:=\sup|c_1|=\frac72r$, $C_2:=\sup|c_2|=\frac52r$, $C_3:=19$, $C_4:=\frac{27}{2}$ (the last is an upper bound, $18\cdot\frac14\cdot3$; true sup $8.93$). Here $r=0.8337$.
Moreover, writing $p_j=\partial_x^{j}h_x$ (so $p_1=h_{xx}$, $p_2=h_{xxx}$, $p_3=h^{(4)}$, $p_4=h^{(5)}$), 
\begin{align}
G_x&=a\,p_3+c_2\,p_1p_2+b'\,p_1^3,\label{eq:Gx}\\
G_{xx}&=a\,p_4+c_1\,p_1p_3+c_2\,p_2^2+c_3\,p_1^2p_2+c_4\,p_1^4 ,\label{eq:Gxx}
\end{align}
with $a,b,c_j$ evaluated at $p=h_x$. (In \eqref{eq:Gx}, \eqref{eq:Gxx} $p_3,p_4$ may involve derivatives of $h$ that are not controlled by (A); they are only ever paired with test functions as shown below.)
\end{lemma}
\begin{proof}
$p\,s^{-3}$ has derivative $(1-5p^2)s^{-4}>0$ on $[0,\frac14]$, so $\sup_{|p|\le1/4}|p|s^{-3}=\frac14r$; this gives $A_1,B_0,C_1,C_2$. For $b'$: $|1-5p^2|\le1$, $s^{-4}\le1$. For $c_3$: $0<1-5p^2\le1$. For $c_4$: $|p|\le\frac14$, $3-5p^2\le3$. \eqref{eq:Gx}--\eqref{eq:Gxx} are obtained by differentiating $G=a(p)p_2+b(p)p_1^2$ with $p=h_x$ (the product/chain rule; e.g.\ the $p_1p_3$ coefficient is $a'+a'+2b$, the $p_1^2p_2$ coefficient is $a''+3b'+2b'$, the $p_2^2$ coefficient is $a'+2b$); the identities and all coefficient formulas are verified symbolically, and the sups numerically, in \texttt{sanity.py}.
\end{proof}

\paragraph{Cutoffs.} Let $S(\sigma)=10\sigma^3-15\sigma^4+6\sigma^5$ on $[0,1]$ ($S=0$ for $\sigma\le0$, $1$ for $\sigma\ge1$). It is $C^2$, nondecreasing, $\max S'=\frac{15}8=:S_1$, $\max|S''|=\frac{10}{\sqrt3}=:S_2$, $\int_0^1S'^2=\frac{10}7$. For a radius $\rho$ and width $\omega$ put $\zeta_{\rho,\omega}(x)=S((\rho+\omega-|x|)/\omega)$: it is even, $C^2$, $=1$ on $|x|\le\rho$, $=0$ for $|x|\ge\rho+\omega$, $|\zeta'|\le S_1/\omega$, $|\zeta''|\le S_2/\omega^2$, $|\{\zeta'\ne0\}|\le2\omega$. Fix $w\in(0,\frac12)$ and use
\[
\zeta_p=\zeta_{1,w},\quad \zeta_q=\zeta_{1+w,w},\quad \zeta=\zeta_E=\zeta_{1+2w,\,w_E},\ w_E=1-2w .
\]
Then $\operatorname{supp}\zeta_p\subset\{\zeta_q=1\}$, $\operatorname{supp}\zeta_q\subset\{\zeta_E=1\}=\{|x|\le r_E\}$, $r_E:=1+2w$. Constants: $Z_p=Z_q=S_1/w$, $Z=S_1/w_E$, $Z_2=S_2/w_E^2$, $m=2w_E$, $q:=12Z^2+4Z_2$. Since $\zeta\in C^2_c(I_0)$ with $\zeta(\pm2)=0$, all integrations by parts below have no boundary terms ($h$ is smooth, so $g_j$ are smooth).

\begin{lemma}[one-dimensional inequalities]\label{lem:1d}
Let $u\in H^1$ have compact support in $I_0$.
\begin{enumerate}
\item[(S)] $\sup u^2\le\nrm u\nrm{u'}$.
\item[(L4)] $\nrm u_{L^4}^4\le\nrm u^3\nrm{u'}$.
\item[(Y)] Let $0\le\zeta\le1$, $|\zeta'|\le Z$, $\zeta\in C^1_c(I_0)$ and $s:=\nrm{\zeta^2g_5}$, $Y:=\nrm{\zeta g_4}$. Then $Y\le 2Z\nrm{g_3}+\sqrt{\nrm{g_3}\,s}$.
\item[(Y$_2$)] Let $\zeta_q$ be as above, $X_1:=\nrm{g_1}_{L^2(|x|\le r_E)}$, $X_3:=\nrm{g_3}_{L^2(|x|\le r_E)}$, $Y_2:=\nrm{\zeta_qg_2}$. Then $Y_2\le 2Z_qX_1+\sqrt{X_1X_3}$.
\end{enumerate}
\end{lemma}
\begin{proof}
(S) $u(x)^2=\int_{-2}^x(u^2)'=-\int_x^2(u^2)'$, so $2u(x)^2\le 2\int|uu'|$, i.e.\ $u^2\le\int|uu'|\le\nrm u\nrm{u'}$. (L4) $\int u^4\le\sup u^2\int u^2$ and (S). (Y) $Y^2=\int\zeta^2g_4^2=-\int g_3(\zeta^2g_4)'=-\int\zeta^2g_3g_5-2\int\zeta\zeta'g_3g_4\le\nrm{g_3}s+2Z\nrm{g_3}Y$; solving the quadratic inequality, $Y\le Z\nrm{g_3}+\sqrt{Z^2\nrm{g_3}^2+\nrm{g_3}s}\le2Z\nrm{g_3}+\sqrt{\nrm{g_3}s}$. (Y$_2$) the same computation with $(\zeta_q,g_1,g_2,g_3)$ in place of $(\zeta,g_3,g_4,g_5)$: $Y_2^2=-\int\zeta_q^2g_1g_3-2\int\zeta_q\zeta_q'g_1g_2\le X_1X_3+2Z_qX_1Y_2$.
\end{proof}

\section{Energy estimate at the level of $g_x$}
\begin{proposition}\label{prop:E1}
Let $E_1(t)=\int\zeta^2g_1^2$ and let $\kappa'\ge\sup_t\nrm{g_3}_{L^2(\operatorname{supp}\zeta)}$ (we use $\kappa'=\kappa$). Then $E_1(t)\le\Lambda_1t$, with
\begin{align*}
\Lambda_1&=\frac{\beta_1^2}{2a_{\min}}+2\gamma_1,\quad \rho_0=M(1+C_2K_2)+B_1K_2^3,\\
\beta_1&=4Zk_2\sqrt m+2A_1K_2k_2+2ZA_1K_2p_1\sqrt m+2C_2K_2k_2+2Zp_1C_2K_2\sqrt m,\\
\gamma_1&=2(Z^2+Z_2)p_1\sqrt m\,\kappa+4k_2\rho_0+2Zp_1\rho_0m .
\end{align*}
\end{proposition}
\begin{proof}
Since $g_t=-G_x$, $\tfrac12E_1'=-\int\zeta^2g_1G_{xx}=\int(\zeta^2g_1)_xG_x$, and by \eqref{eq:Gx} $G_x=ag_4+r_1$, $r_1=ah_0^{(4)}+c_2h_{xx}h_{xxx}+b'h_{xx}^3$. With $(\zeta^2g_1)_x=\zeta^2g_2+2\zeta\zeta'g_1$:
\begin{align*}
\tfrac12E_1'&=\underbrace{\int\zeta^2ag_2g_4}_{T_1}+\underbrace{\int\zeta^2g_2r_1}_{T_2}+\underbrace{2\int\zeta\zeta'ag_1g_4}_{T_3}+\underbrace{2\int\zeta\zeta'g_1r_1}_{T_4}.
\end{align*}
Put $\sigma:=\nrm{\zeta g_3}$; then $\int\zeta^2|g_3|\le\nrm\zeta\sigma\le2\sigma$ and $\int_{\{\zeta'\ne0\}}\zeta|g_3|\le\sqrt m\,\sigma$. Integrating by parts once ($a_x=a'h_{xx}$):
\begin{align*}
T_1&=-\int\zeta^2ag_3^2-\int(2\zeta\zeta'a+\zeta^2a_x)g_2g_3\ \le\ -a_{\min}\sigma^2+2Zk_2\sqrt m\,\sigma+2A_1K_2k_2\,\sigma,\\
T_3&=-2\int\big[(\zeta'^2+\zeta\zeta'')ag_1+\zeta\zeta'a_xg_1+\zeta\zeta'ag_2\big]g_3\ \le\ 2(Z^2+Z_2)p_1\sqrt m\,\kappa'+2ZA_1K_2p_1\sqrt m\,\sigma+2Zk_2\sqrt m\,\sigma .
\end{align*}
Pointwise, $|r_1|\le\rho_0+C_2K_2|g_3|$ (using $|h_{xxx}|\le M+|g_3|$), so
$T_2\le4k_2\rho_0+2k_2C_2K_2\sigma$ and $T_4\le2Zp_1\rho_0m+2Zp_1C_2K_2\sqrt m\,\sigma$. Collecting, $\frac12E_1'\le-a_{\min}\sigma^2+\beta_1\sigma+\gamma_1\le\beta_1^2/(4a_{\min})+\gamma_1$. Since $E_1(0)=0$, $E_1(t)\le\Lambda_1t$.
\end{proof}

\section{Energy estimate at the level of $g_{xxx}$}
\begin{proposition}\label{prop:E3}
Let $E_3(t)=\int\zeta^4g_3^2$, $K:=\kappa$ ($\ge\sup_t\nrm{g_3}_{L^2(\operatorname{supp}\zeta)}$). Define the polynomials in $u\ge0$
\begin{align*}
\bar Y(u)&=2ZK+K^{1/2}u^2,\qquad \bar\omega(u)=K^{3/2}\big[(3ZK)^{1/2}+K^{1/4}u\big],\\
\bar\nu(u)&=\nu_0+C_1K_2\bar Y(u)+2C_2\bar\omega(u),\qquad \nu_0=M+2C_1K_2M+4C_2M^2+C_3K_2^2(2M+K)+2C_4K_2^4,\\
f(u)&=-a_{\min}u^8+\big(8Z\bar Y(u)+qK\big)u^4+\big(u^4+8Z\bar Y(u)+qK\big)\bar\nu(u).
\end{align*}
Then $E_3(t)\le C_*t$ with $C_*:=2\max_{u\ge0}f(u)$.
\end{proposition}
\begin{proof}
\emph{Step 1 (weak form).} $\partial_tg_3=-\partial_x^4G$, so integrating by parts twice, $\tfrac12E_3'=-\int(\zeta^4g_3)''\,G_{xx}$, and $(\zeta^4g_3)''=\zeta^4g_5+8\zeta^3\zeta'g_4+Qg_3$, $Q=12\zeta^2\zeta'^2+4\zeta^3\zeta''$, $|Q|\le q\zeta^2$.

\emph{Step 2 (splitting).} By \eqref{eq:Gxx} with $p_4=h_0^{(5)}+g_5$, $p_3=h_0^{(4)}+g_4$, $p_2=h_0'''+g_3$: $G_{xx}=ag_5+N$,
\[
N=ah_0^{(5)}+c_1h_{xx}(h_0^{(4)}+g_4)+c_2h_{xxx}^2+c_3h_{xx}^2h_{xxx}+c_4h_{xx}^4 .
\]
Let $s:=\nrm{\zeta^2g_5}$, $Y:=\nrm{\zeta g_4}$, $\nu:=\nrm{\zeta^2N}$. Using $\zeta^4g_5N=(\zeta^2g_5)(\zeta^2N)$, $\zeta^3\zeta'g_4=(\zeta' )(\zeta g_4)\zeta^2$, $|Q g_3|\le q\zeta^2|g_3|$ and Cauchy--Schwarz,
\[
\tfrac12E_3'\le-a_{\min}s^2+(8ZY+qK)s+(s+8ZY+qK)\nu .
\]
(The six terms are: main dissipation $-\int a\zeta^4g_5^2$; $\int\zeta^4g_5N\le s\nu$; $8\int\zeta^3\zeta'ag_4g_5\le8ZYs$; $8\int\zeta^3\zeta'g_4N\le8ZY\nu$; $\int Qag_3g_5\le qKs$; $\int Qg_3N\le qK\nu$.)

\emph{Step 3 (bounds for $Y,\nu$ in terms of $s$).} Write $N=N_1+\dots+N_5$ in the order displayed. By (Y), $Y\le\bar Y(s^{1/4})$. For $\nu$: $\nrm{\zeta^2ah_0^{(5)}}\le M$; $\nrm{\zeta^2N_2}\le C_1K_2(2M+Y)$ (as $\nrm{\zeta^2}\le2$, $|h_0^{(4)}|\le M$); $\nrm{\zeta^2N_4}\le C_3K_2^2(2M+K)$; $\nrm{\zeta^2N_5}\le2C_4K_2^4$; and, by (L4) applied to $u=\zeta g_3$ (with $\nrm{u'}\le Y+ZK$),
$\nrm{\zeta^2N_3}\le C_2\nrm{\zeta h_{xxx}}_{L^4}^2\le C_2(4^{1/4}M+\omega)^2\le C_2(4M^2+2\omega^2)$, $\omega^2:=\nrm{\zeta g_3}_{L^4}^2\le K^{3/2}(Y+ZK)^{1/2}\le\bar\omega(s^{1/4})$ (using $\sqrt{x+y}\le\sqrt x+\sqrt y$). Hence $\nu\le\bar\nu(s^{1/4})$; all of $\bar Y,\bar\nu$ have nonnegative coefficients, so the bound is monotone in $s$.

\emph{Step 4.} With $u=s^{1/4}$: $\frac12E_3'\le f(u)\le\max f$. $f$ is a polynomial of degree 8 with leading coefficient $-a_{\min}<0$, so the maximum exists and is attained at $u=0$ or a positive root of $f'$ (a degree-7 polynomial); \texttt{lemmaO\_consts.py} computes it from these roots (and cross-checks on a grid). $E_3(0)=0$ gives the claim.
\end{proof}

\section{Mean bound and interpolation; proof of Theorem \ref{thm:O}}
\begin{lemma}[weak bound]\label{lem:mean}
Let $\theta=\zeta_p$. Then $\big|\int\theta\,g(\cdot,t)\big|\le t\,\nrm{\theta'}\,\Gamma$, with $\nrm{\theta'}=\sqrt{2S_{11}/w}$, $S_{11}=\frac{10}7$.
\end{lemma}
\begin{proof}
$\frac d{dt}\int\theta g=-\int\theta G_x=\int\theta'G(h)$, $|\cdot|\le\nrm{\theta'}\Gamma$ by \eqref{eq:P}; $g(\cdot,0)=0$.
\end{proof}

\begin{proof}[Proof of Theorem \ref{thm:O}]
Choose $w$ (the script scans $w\in[0.04,0.46]$ and takes the best; any $w$ is valid). By Props.\ \ref{prop:E1}, \ref{prop:E3}, since $\zeta_E=1$ on $|x|\le r_E$, for $t\le t_1$
\[
X_1^2\le E_1\le\Lambda_1t,\qquad X_3^2\le E_3\le C_*t,\qquad\text{in particular }\int_{-1}^1g_3^2\le C_*t\ \ (C_3:=C_*).
\]
\emph{$g_{xx}$ in $L^2$.} By (Y$_2$), $Y_2=\nrm{\zeta_qg_2}\le Y_c\sqrt t$, $Y_c:=2Z_q\sqrt{\Lambda_1}+(\Lambda_1C_*)^{1/4}$.
\emph{$g_x$ pointwise.} Apply (S) to $u=\zeta_qg_1$ ($\zeta_q=1$ on $|x|\le1+w$): $\sup_{|x|\le1}g_1^2\le X_1(Y_2+Z_qX_1)\le\sqrt{\Lambda_1}(Y_c+Z_q\sqrt{\Lambda_1})\,t$.
\emph{$g_{xx}$ pointwise.} Apply (S) to $u=\zeta_pg_2$: $\nrm u\le Y_2$ (as $\zeta_p\le\zeta_q=1$ on $\operatorname{supp}\zeta_p$), $\nrm{u'}\le\nrm{\zeta_pg_3}+Z_p\nrm{g_2}_{L^2(\operatorname{supp}\zeta_p)}\le X_3+Z_pY_2$. So $\sup_{|x|\le1}g_2^2\le Y_c(\sqrt{C_*}+Z_pY_c)\,t$.
\emph{$g$ pointwise.} For $|x|\le1$ and $y\in\operatorname{supp}\theta\subset[-1-w,1+w]$, $|g(x)-g(y)|\le|x-y|^{1/2}\nrm{g_1}_{L^2(|z|\le 1+w)}\le\sqrt{2+w}\sqrt{E_1}$. As $\theta\in[0,1]$, $\int\theta\ge2$,
\[
|g(x,t)|\le\frac{|\int\theta g|}{\int\theta}+\sqrt{2+w}\sqrt{\Lambda_1t}\le \tfrac12\nrm{\theta'}\Gamma\,t+\sqrt{2+w}\sqrt{\Lambda_1}\,t^{1/2},
\]
and $t\le\sqrt{t_1}\,t^{1/2}$. Altogether
\begin{equation}\label{eq:consts}
\boxed{\begin{aligned}
C_0&=\sqrt{2+w}\sqrt{\Lambda_1}+\tfrac12\sqrt{2S_{11}/w}\,\Gamma\sqrt{t_1},\\
C_1&=\big[\sqrt{\Lambda_1}(Y_c+Z_q\sqrt{\Lambda_1})\big]^{1/2},\quad
C_2=\big[Y_c(\sqrt{C_*}+Z_pY_c)\big]^{1/2},\quad C_3=C_*,
\end{aligned}}
\end{equation}
with $\Lambda_1,C_*$ from Props.\ \ref{prop:E1}, \ref{prop:E3} and $Y_c$ above.
\end{proof}

\begin{remark}[$W^{2,\infty}$ cutoffs]
The cutoff $\zeta_E$ is $C^2$ but not $C^\infty$. The user-requested $\eta\in C_c^\infty$ is not needed: every integration by parts above involves at most two derivatives of $\zeta^4$, resp.\ $\zeta^2$, which are continuous.
\end{remark}

\section{Numerical values and the two use-cases}\label{sec:use}
\paragraph{What is computed.} \texttt{lemmaO\_consts.py} evaluates \eqref{eq:consts}. It is double-precision arithmetic, not interval arithmetic; the polynomial maximum in Prop.~\ref{prop:E3} is inflated by a relative $10^{-9}$. Every formula in the script is the one in the Props.\ above.

\begin{table}[h]\centering\small
\begin{tabular}{ccccccccc}\toprule
$(K_2,K_3,M)$ & $t_1$ & $C_0$ & $C_1$ & $C_2$ & $C_3$ & $C_1\sqrt{t_1}$ & $C_2\sqrt{t_1}$ & $\sqrt{C_3t_1}$\\\midrule
$(0.1,1,0.1)$ & $3.2\cdot10^{-4}$ & 10.1 & 46.3 & 916 & $6.8\cdot10^{5}$ & 0.83 & 16 & 15\\
$(0.3,1,0.3)$ & $3.2\cdot10^{-4}$ & 20.8 & 86.9 & $1.5\cdot10^3$ & $2.3\cdot10^6$ & 1.6 & 26 & 27\\
$(1,1,1)$ & $3.2\cdot10^{-4}$ & 86.6 & 330 & $4.8\cdot10^3$ & $2.3\cdot10^7$ & 5.9 & 85 & 86\\
$(10,10,10)$ & $10^{-8}$ & $3.9\cdot10^3$ & $2.0\cdot10^4$ & $4.7\cdot10^5$ & $2.8\cdot10^{11}$ & 2.0 & 47 & 53\\\bottomrule
\end{tabular}\end{table}
Trivial sizes from (A),(I): $|g_x|\le0.25+M$, $|g_{xx}|\le K_2+M$, $\nrm{g_3}\le K_3+2M$.

\paragraph{Verdict.} The theorem is true for all $t_1$, but for both use-cases the bounds are \emph{vacuous}: $C_2\sqrt{t_1}$ is 16--85 (use-case (i)) and 47 (use-case (ii)) against trivial sizes $\lesssim 2$ and $20$. Requiring all four deviations $\le0.1$ needs $t\lesssim 1.2\cdot10^{-8}$ ($K_2=M=0.1$), $4\cdot10^{-10}$ ($K_2=M=1$) in case (i) --- i.e.\ $t_1$ smaller by factors $\approx3\cdot10^4$ to $7\cdot10^5$ than $3.2\cdot10^{-4}$ --- and $t\lesssim3.5\cdot10^{-14}$ in case (ii) (factor $3\cdot10^5$ below $10^{-8}$). In terms of $\delta_0$ for ``$C_j\sqrt{t_1}\le0.1$'': $\delta_0\approx1.2\cdot10^{-8}$, $4.4\cdot10^{-9}$, $4.4\cdot10^{-10}$ for $K_2=M=0.1,0.3,1$ ($K_3=1$), and $3.5\cdot10^{-14}$ for $K_2=K_3=M=10$.

\paragraph{Bottlenecks (diagnosed with the script).} (1) $C_3$: the dominant factor is $\kappa=K_3+2M$, the \emph{unweighted} $L^2$ norm of $g_{xxx}$ on the cutoff transition region, multiplied by the cutoff curvature constant $q=12Z^2+4Z_2$ (terms $qKs$ and $qK\nu$ in Prop.~\ref{prop:E3}). Replacing $\kappa$ by $\kappa/10$ lowers $C_3$ at $(1,1,1)$ from $2.3\cdot10^7$ to $3\cdot10^5$ (factor 75), and $\kappa\to0.01$ gives $1.2\cdot10^4$, which is the floor set by the nonlinear coefficients ($C_3K_2^2$, $C_4K_2^4$ in $\nu_0$). Dividing $q$ by 10 helps by only a factor $\approx1.4$. (2) $\Lambda_1$ (hence $C_1,C_2$): the term $\beta_1^2/(2a_{\min})$ with $4Zk_2\sqrt m$ (cutoff gradient times the unweighted pointwise bound $k_2=K_2+M$ of $g_{xx}$) and $\gamma_1$'s $2(Z^2+Z_2)p_1\sqrt m\,\kappa$. (3) $C_2$ and $C_0$ inherit these via $Y_c\sim2Z_q\sqrt{\Lambda_1}+(\Lambda_1C_*)^{1/4}$; $C_2\approx\sqrt{Y_c\sqrt{C_*}}$ is limited by $C_*$.

\paragraph{Heuristic remedy (not proved here).} Replace $\kappa$ by the \emph{small} quantity $\sqrt{C_*^{\rm outer}t}$ using a chain of nested cutoffs (the transition region of an inner cutoff lies in the plateau of an outer one, where $\int g_3^2\le C_*^{\rm outer}t$); similarly for $k_2,p_1$ via (Y$_2$). A trial of this for $E_3$ alone (code removed because it is not part of the proof) reduced $C_3$ for $(1,1,1)$ by a factor $\approx7$ at $t_1=10^{-9}$ (two levels) and by a factor $\approx10^{3}$ at $t_1=10^{-12}$ (three or four levels), but it does not help at $t_1\gtrsim10^{-7}$, because $C_*t_1$ must be below $\kappa^2$ at the outermost level, and it does not reduce $\Lambda_1$ unless $E_1$ is chained too.

\appendix
\section{Every inequality used, with its constant}
\begin{center}\small\begin{tabular}{p{0.08\textwidth}p{0.62\textwidth}p{0.24\textwidth}}\toprule
 & Inequality & Constant\\\midrule
I1 & $\Amin\le a\le1$, $|a'|\le A_1$, $|b|\le B_0$, $|b'|\le B_1$, $|c_{1,2,3,4}|\le C_{1,2,3,4}$ on $|p|\le\frac14$ (Lemma \ref{lem:coef}) & $0.8858$; $0.8337$; $0.6253$; $3$; $2.918,\,2.084,\,19,\,13.5$\\
I2 & Quintic smoothstep: $\max S'$, $\max|S''|$, $\int S'^2$ & $15/8$; $10/\sqrt3$; $10/7$\\
I3 & Pointwise/integral consequences of (A),(I): $|g_1|\le p_1$, $|g_2|\le k_2$, $\nrm{g_3}\le\kappa$, $\nrm{G}\le\Gamma$ & $\frac14+M$; $K_2+M$; $K_3+2M$; $K_3+2B_0K_2^2$\\
I4 & $\nrm{\zeta}_{L^2(I_0)},\nrm{\zeta^2}_{L^2(I_0)}\le2$; $\int_{\{\zeta'\ne0\}}\zeta|g_3|\le\sqrt m\,\nrm{\zeta g_3}$, $m=2w_E$ & $2$; $\sqrt{2w_E}$\\
I5 & Cauchy--Schwarz in all pairings of Props.\ \ref{prop:E1}, \ref{prop:E3} & 1\\
I6 & $\sup(-\Amin\sigma^2+\beta\sigma)=\beta^2/(4\Amin)$ (Prop.~\ref{prop:E1}) & --\\
I7 & (S) $\sup u^2\le\nrm u\nrm{u'}$ (compact support) & 1\\
I8 & (L4) $\nrm u_4^4\le\nrm u^3\nrm{u'}$ & 1\\
I9 & (Y) $Y\le2Z\nrm{g_3}+\sqrt{\nrm{g_3}s}$; (Y$_2$) $Y_2\le2Z_qX_1+\sqrt{X_1X_3}$ & 2\\
I10 & $(x+y)^2\le2x^2+2y^2$; $\sqrt{x+y}\le\sqrt x+\sqrt y$; $\nrm{\zeta}_{L^4}\le4^{1/4}$ & $2$; $1$; $4^{1/4}$\\
I11 & Max of the degree-8 polynomial $f$ (positive-coefficient perturbation of $-\Amin u^8$) via roots of $f'$ & numerical, $+10^{-9}$ rel.\\
I12 & Mean value: $|g(x)-g(y)|\le|x-y|^{1/2}\nrm{g_1}$, $|x-y|\le2+w$; $\int\theta\ge2$; $t\le\sqrt{t_1}\sqrt t$ & $\sqrt{2+w}$\\\bottomrule
\end{tabular}\end{center}

\paragraph{What is not interval-certified.} (a) Floating-point evaluation of \eqref{eq:consts}; (b) the polynomial maximum of I11 (root finding); (c) the closed-form sups in I1 are confirmed on a fine grid in \texttt{sanity.py}; the true $\sup|c_4|=8.93$ is smaller than the $C_4=13.5$ used, so that bound is safe. Everything else is a closed-form inequality chain.
\end{document}
